\newpage
\begin{center}
    {\large \textbf{Olivaris - Plataforma de gestión colaborativa para el olivar}}

    \vspace{0.2cm}
    {\normalsize Jorge Cano Melero}

    \vspace{0.2cm}
\end{center}

\vspace{0.5cm}

\textbf{Keywords:} Commercial Farm Logbook, Agricultural Management, Backend, REST API, JWT, SIGPAC, Parcels, Phytosanitary Treatments.

\vspace{0.5cm}

\textbf{Abstract:}

This Final Degree Project addresses the design and implementation of Olivaris, a digital platform aimed at collaborative management of olive grove farms. The project arises in a context of increasing digitalization of 
the agricultural sector and regulatory evolution around the System of Information of Agricultural Farms, the Regional Registry of Agricultural Farms, and the Digital Farm Operation Logbook. Within this framework, Olivaris 
is presented as a tool to centralize the management of parcels, plots, users, authorized entities and agricultural activities, especially those related to the use of phytosanitary products.

The developed solution is based on a client-server architecture, with a backend that exposes a REST API, database persistence, authentication through JWT and role-based access control. The system allows registering and viewing 
parcels, managing activities associated with farms and modeling the relationship between farmers and authorized entities through access and management permissions.

Although the platform does not perform an effective integration with the official SIEX/REA services, its design has been developed taking into account the conceptual structure and information requirements associated with this 
framework, in order to facilitate a possible future integration. As a result, a functional and extensible prototype is obtained that contributes to the digitalization of olive grove management and serves as a basis for future 
improvements, such as offline functionality, advanced report generation or cross-platform access.